\honest{Expecting Short Inferential Distances}{Eliezer Yudkowsky}{2007}

Our ancestors lived in bands of ``at most 200 people,'' with no writing. ``In a world like that, all background knowledge is universal knowledge. All information not strictly private is public, period.'' The word ``period'' is the whole argument for this. Every band had children, who knew none of what the adults knew and had to be taught all of it, over years. I leave the children out.

The rest follows from that premise. If you found an oasis, you never had to explain what an oasis was. No one's conclusions were ``a hundred inferential steps removed'' from what everyone knew. So a person who said something with no obvious support was ``a liar or an idiot,'' and a person who missed something obvious was an idiot too, or was annoying you on purpose.

This, I think, explains ``a lot'' about why scientists explain things badly to the public and to scientists in other fields. Both sides expect new ideas to be only a short distance from what everyone knows. In the same sentence I name two other effects that combine with this one: the illusion of transparency, and the habit of modelling other minds as slightly modified copies of your own. Either of those alone would produce the failures I describe. The story about hunter-gatherers is a third cause, and nothing in the post would come out differently if it were false. My evidence that the failures happen is what I see ``When I observe failures of explanation.'' Economists had named the effect ``the curse of knowledge'' in 1989. I do not mention them.

My example: a biologist can justify evolution to a physicist by calling it ``the simplest explanation.'' Someone not raised on the history of science from Newton to Einstein hears an interesting point, not a knockdown argument, and decides the biologist is arrogant. The biologist decides that ``nonscientists are just idiots.''

Then the advice. Build a path of arguments that starts from what the audience ``already knows or accepts''; ``If you don't recurse far enough, you're just talking to yourself.'' Do not visibly give an argument more weight than the audience does, as when you treat ``simpler explanation'' as a knockdown argument for evolution ``(which it is).'' And do not hint that you are a dozen steps ahead of them, or ``They'll just think you're condescending.''

I end: if you think you can explain systematically underestimated inferential distances ``briefly, in just a few words, I've got some sad news for you.'' I have just explained it in about 800 words.
